\bookchapter{第三十一章：短影片拍攝與內容創作入門}

\chapterintro{2026年，短影片已經唔係「後生仔嘅玩具」喇，佢變成所有人獲取資訊同表達自我嘅基礎方式。唔管你想記錄生活、分享知識，仲係做啲副業，手機就係你最好嘅創作工具——唔使買相機、唔使學Premiere，一部手機就能開始。呢一章幫你由零起步。}

\booksection{一、拍之前先諗清楚：你要拍畀邊個睇？}

好多人拍影片嘅第一個錯誤就係「攞起手機就拍」，拍完唔知畀邊個睇。喺開拍之前，先回答三個問題：

\begin{enumerate}
  \item \textbf{我嘅觀眾係邊個？}——拍畀仔睇、拍畀同齡人睇、拍畀長輩睇，內容風格完全唔同。
  \item \textbf{我能提供乜嘢價值？}——等人笑（娛樂）、等學嘢（知識）、等人饞（美食）、等人暖（情感）。至少佔一樣。
  \item \textbf{我能喺邊個平臺持續更新？}——抖音/快手（泛娛樂）、小紅書（生活方式）、B站（深度內容）、影片號（社交圈）。選一個主陣地深耕，唔好一嚟就全平臺分發。
\end{enumerate}

\booksection{二、手機拍攝基礎設定——開拍前必做嘅5件事}

\begin{enumerate}
  \item \textbf{解像度設1080P，幀率設30fps}：呢個係短影片平臺嘅標準。4K檔案太大上傳慢，60fps喺社交平臺上冇優勢反而檔案更大。
  \item \textbf{抹乾淨鏡頭}：第一條亦係最重要嘅一條。拍之前用衫抹一下鏡頭，畫面立刻清晰一個檔次。
  \item \textbf{鎖定對焦同曝光}：喺取景框長按畫面主體2秒，會出現「AE/AF鎖定」提示。咁樣拍嘅時候唔會因為移動導致畫面忽明忽暗。
  \item \textbf{橫拍定豎拍？}——發抖音/快手/影片號/小紅書\textbf{豎住拍}（9:16），發B站/YouTube\textbf{橫住拍}（16:9）。諗清楚平臺再決定方向。
  \item \textbf{開網格線}：設定 > 相機 > 網格線。三分法構圖嘅基礎，幫你快速搵到畫面重心。
\end{enumerate}

\booksection{三、運鏡五招——等畫面「鬱」起嚟}

新手拍影片最大嘅問題係畫面太「靜」——全程企住唔鬱拍。學識下面五個基本運鏡，畫面立刻有專業感：

\begin{itemize}
  \item \textbf{推（Push）}：手機慢慢向被拍對象靠近。製造「靠近」嘅緊張感或親密感。
  \item \textbf{拉（Pull）}：手機慢慢遠離被拍對象。由細節拉到全景，交代環境。
  \item \textbf{搖（Pan）}：手機位置唔鬱，左右轉動。適合影風景、展示空間。
  \item \textbf{移（Track）}：手機平行移動，跟住被拍對象行。需要手持穩定，最好配合手機穩定器。
  \item \textbf{跟（Follow）}：手機跟喺被拍對象後面或前面行。Vlog最常用嘅運鏡。
\end{itemize}

\textbf{穩係第一要務}：手震嘅畫面睇落好唔專業。雙手持機比單手穩；行路時膝頭微屈，用「忍者步」減少上下顛簸；如果有預算，買一個手機穩定器（大疆OM系列約500-800蚊），效果質變。

\booksection{四、剪輯入門——剪映就夠}

剪映（CapCut嘅國內版）係目前最適合新手嘅手機剪輯工具，免費、功能全、模板多。以下係基礎工作流：

\begin{enumerate}
  \item \textbf{匯入素材}：打開剪映 > 開始創作 > 選擇拍好嘅影片片段匯入。
  \item \textbf{裁剪}：拖動時間軸，將每個片段開頭結尾嘅廢片剪掉。原則：每一段淨係保留有資訊量嘅部分，寧短勿長。
  \item \textbf{加轉場}：兩段影片之間點「轉場」按鈕，選一個簡單嘅轉場效果。新手建議用「疊化」或「黑場」，唔好用花哨嘅3D轉場。
  \item \textbf{加字幕}：剪映有「識別字幕」功能，一鍵自動生成字幕。準確率約90\%，手動校對一遍即可。
  \item \textbf{配音樂}：剪映音樂庫入面便有大量版權免費嘅音樂。選一首節奏同內容匹配嘅BGM，音量調到影片原聲嘅1/3左右。
  \item \textbf{匯出}：1080P匯出，幀率30fps，碼率選「推薦」即可。
\end{enumerate}

\booksection{五、封面嘅重要性——70\%嘅人淨係睇封面決定睇唔睇}

短影片嘅封面就係「門面」，決定別人要不要點入嚟。幾個做封面嘅要點：

\begin{itemize}
  \item \textbf{封面要有一句話}：用大字寫明影片嘅核心內容，比如「3分鐘學識手機修圖」。字體要大、顏色要亮、同背景有對比。
  \item \textbf{封面同內容要一致}：唔好做「標題黨」——封面講一件事，點入去係另一回事，觀眾會立刻劃走，算法會降權。
  \item \textbf{人物表情要生動}：如果封面有人，選一個表情生動或誇張嘅瞬間。人臉比文字更能吸引㩒。
  \item \textbf{剪映做封面}：喺剪映入面便選一幀最好嘅畫面，加文字疊加，匯出為圖片作為封面。
\end{itemize}

\booksection{六、內容創作嘅三條心法}

\begin{itemize}
  \item \textbf{前3秒定生死}：短影片平臺，觀眾3秒內決定掃走定係繼續睇。前3秒必須畀出「鈎子」——一個問題、一個懸念、一個意外畫面。
  \item \textbf{一條影片淨係講一件事}：新手最容易犯嘅錯係一條影片塞太多內容。一條影片淨係解決一個問題，講透比講全重要。
  \item \textbf{堅持比天賦重要}：算法推薦需要「養號」。前20條影片幾乎唔會有播放量，呢個係正常嘅。堅持每週更新2-3條，等算法認識你嘅內容方向，流量會慢慢起嚟。
\end{itemize}

\begin{tcbpractice}{動手：拍一條30秒嘅短影片}
\begin{enumerate}
  \item 選一個你擅長嘅小知識（比如「點樣將手機字體調大」）。
  \item 用手機豎屏拍攝3-5段素材，每段5-10秒。
  \item 用剪映接駁、裁剪、加字幕、配BGM。
  \item 做一張帶大字嘅封面圖。
  \item 匯出後先發畀朋友睇，收集反饋再發布。
\end{enumerate}

\end{tcbpractice}

\begin{tcbsummary}{本章小結}
\begin{itemize}
  \item 拍之前諗清楚：觀眾係邊個、提供乜價值、喺邊個平臺發
  \item 開拍5件事：1080P/30fps、抹鏡頭、鎖對焦曝光、確定橫豎、開網格線
  \item 運鏡五招：推、拉、搖、移、跟——穩係第一要務
  \item 剪輯用剪映：裁剪轉場字幕BGM匯出，六步走
  \item 前3秒定生死，一條影片淨係講一件事，堅持比天賦重要
\end{itemize}

\end{tcbsummary}
